\documentclass{report}
\usepackage{luatexja} % LuaTeXで日本語を使うためのパッケージ
\usepackage{luatexja-fontspec} % LuaTeX用の日本語フォント設定

% --- 数学関連 ---
\usepackage{amsmath, amssymb, amsfonts, mathtools, bm, amsthm} % 基本的な数学パッケージ
\usepackage{type1cm, upgreek} % 数式フォントとギリシャ文字
\usepackage{physics, mhchem} % 物理や化学の記号や式の表記を簡単にする

% --- 表関連 ---
\usepackage{multirow, longtable, tabularx, array, colortbl, dcolumn, diagbox} % 表のレイアウトを柔軟にする
\usepackage{tablefootnote, truthtable} % 表中に注釈を追加、真理値表
\usepackage{tabularray} % 高度な表組みレイアウト

% --- グラフィック関連 ---
\usepackage{tikz, graphicx} % 図の描画と画像の挿入
\usepackage{background} % ウォーターマークの設定
\usepackage{caption, subcaption} % 図や表のキャプション設定
\usepackage{float, here} % 図や表の位置指定

% --- レイアウトとページ設定 ---
\usepackage{fancyhdr} % ページヘッダー、フッター、余白の設定
\usepackage[top = 20truemm, bottom = 20truemm, left = 20truemm, right = 20truemm]{geometry}
\usepackage{fancybox, ascmac} % ボックスのデザイン

% --- 色とスタイル ---
\usepackage{xcolor, color, colortbl, tcolorbox} % 色とカラーボックス
\usepackage{listings, jvlisting} % コードの色付けとフォーマット

% --- 参考文献関連 ---
\usepackage{biblatex, usebib} % 参考文献の管理と挿入
\usepackage{url, hyperref} % URLとリンクの設定

% --- その他の便利なパッケージ ---
\usepackage{footmisc} % 脚注のカスタマイズ
\usepackage{multicol} % 複数段組
\usepackage{comment} % コメントアウトの拡張
\usepackage{siunitx} % 単位の表記
\usepackage{docmute}
% \usepackage{appendix}
% --- tcolorboxとtikzの設定 ---
\tcbuselibrary{theorems, breakable} % 定理のボックスと改ページ設定
\usetikzlibrary{decorations.markings, arrows.meta, calc} % tikzの装飾や矢印の設定

% --- 定理スタイルと数式設定 ---
\theoremstyle{definition} % 定義スタイル
\numberwithin{equation}{section} % 式番号をサブセクション単位でリセット

% --- hyperrefの設定 ---
\hypersetup{
  setpagesize = false,
  bookmarks = true,
  bookmarksdepth = tocdepth,
  bookmarksnumbered = true,
  colorlinks = false,
  pdftitle = {}, % PDFタイトル
  pdfsubject = {}, % PDFサブジェクト
  pdfauthor = {}, % PDF作者
  pdfkeywords = {} % PDFキーワード
}

% --- siunitxの設定 ---
\sisetup{
  table-format = 1.5, % 小数点以下の桁数
  table-number-alignment = center, % 数値の中央揃え
}

% --- 透かし画像の設定 --- 
\backgroundsetup{
  scale=0.5,                       % 画像のスケール
  % color=black,                   % 画像の色（透かし用に半透明が推奨）
  opacity=0.2,                   % 透かしの透明度（0が完全透明、1が完全不透明）
  angle=0,                       % 画像の角度
  position = current page.south east,  % ページの右下
  hshift=-6cm, % 右方向へのシフト（負の値で内側に移動）
  vshift=5cm,  % 上方向へのシフト（正の値で内側に移動）
  contents={\includegraphics{./fig/appilogo-circular-full.png}} % 画像のパス
}

% --- その他の設定 ---
\allowdisplaybreaks % 数式の途中改ページ許可
\newcolumntype{t}{!{\vrule width 0.1pt}} % 新しいカラムタイプ
\newcolumntype{b}{!{\vrule width 1.5pt}} % 太いカラム
\UseTblrLibrary{amsmath, booktabs, counter, diagbox, functional, hook, html, nameref, siunitx, varwidth, zref} % tabularrayのライブラリ
\setlength{\columnseprule}{0.4pt} % カラム区切り線の太さ
\captionsetup[figure]{font = bf} % 図のキャプションの太字設定
\captionsetup[table]{font = bf} % 表のキャプションの太字設定
\captionsetup[lstlisting]{font = bf} % コードのキャプションの太字設定
\captionsetup[subfigure]{font = bf, labelformat = simple} % サブ図のキャプション設定
\setcounter{secnumdepth}{4} % セクションの深さ設定
\newcolumntype{d}{D{.}{.}{5}} % 数値のカラム
\newcolumntype{M}[1]{>{\centering\arraybackslash}m{#1}} % センター揃えのカラム
\DeclareMathOperator{\diag}{diag}
\everymath{\displaystyle} % 数式のスタイル
\newcommand{\inner}[2]{\left\langle #1, #2 \right\rangle}
\renewcommand{\figurename}{図}
\renewcommand{\i}{\mathrm{i}} % 複素数単位i
\renewcommand{\laplacian}{\grad^2} % ラプラシアンの記号
\renewcommand{\thesubfigure}{(\alph{subfigure})} % サブ図の番号形式
\newcommand{\m}[3]{\multicolumn{#1}{#2}{#3}} % マルチカラムのショートカット
\renewcommand{\r}[1]{\mathrm{#1}} % mathrmのショートカット
\newcommand{\e}{\mathrm{e}} % 自然対数の底e
\newcommand{\Ef}{E_{\mathrm{F}}} % フェルミエネルギー
\renewcommand{\c}{\si{\degreeCelsius}} % 摂氏記号
\renewcommand{\d}{\r{d}} % d記号
\renewcommand{\t}[1]{\texttt{#1}} % タイプライタフォント
\newcommand{\kb}{k_{\mathrm{B}}} % ボルツマン定数
% \renewcommand{\phi}{\varphi} % ϕをφに変更
\renewcommand{\epsilon}{\varepsilon}
\newcommand{\fullref}[1]{\textbf{\ref{#1} \nameref{#1}}}
\newcommand{\reff}[1]{\textbf{図\ref{#1}}} % 図参照のショートカット
\newcommand{\reft}[1]{\textbf{表\ref{#1}}} % 表参照のショートカット
\newcommand{\refe}[1]{\textbf{式\eqref{#1}}} % 式参照のショートカット
\newcommand{\refp}[1]{\textbf{コード\ref{#1}}} % コード参照のショートカット
\renewcommand{\lstlistingname}{コード} % コードリストの名前
\renewcommand{\theequation}{\thesection.\arabic{equation}} % 式番号の形式
\renewcommand{\footrulewidth}{0.4pt} % フッターの線
\newcommand{\mar}[1]{\textcircled{\scriptsize #1}} % 丸囲み文字
\newcommand{\combination}[2]{{}_{#1} \mathrm{C}_{#2}} % 組み合わせ
\newcommand{\thline}{\noalign{\hrule height 0.1pt}} % 細い横線
\newcommand{\bhline}{\noalign{\hrule height 1.5pt}} % 太い横線

% --- カスタム色定義 ---
\definecolor{burgundy}{rgb}{0.5, 0.0, 0.13} % バーガンディ色
\definecolor{charcoal}{rgb}{0.21, 0.27, 0.31} % チャコール色
\definecolor{forest}{rgb}{0.0, 0.35, 0} % 森の緑色

% --- カスタム定理環境の定義 ---
\newtcbtheorem[number within = chapter]{myexc}{練習問題}{
  fonttitle = \gtfamily\sffamily\bfseries\upshape,
  colframe = forest,
  colback = forest!2!white,
  rightrule = 1pt,
  leftrule = 1pt,
  bottomrule = 2pt,
  colbacktitle = forest,
  theorem style = standard,
  breakable,
  arc = 0pt,
}{exc-ref}
\newtcbtheorem[number within = chapter]{myprop}{命題}{
  fonttitle = \gtfamily\sffamily\bfseries\upshape,
  colframe = blue!50!black,
  colback = blue!50!black!2!white,
  rightrule = 1pt,
  leftrule = 1pt,
  bottomrule = 2pt,
  colbacktitle = blue!50!black,
  theorem style = standard,
  breakable,
  arc = 0pt
}{proposition-ref}
\newtcbtheorem[number within = chapter]{myrem}{注意}{
  fonttitle = \gtfamily\sffamily\bfseries\upshape,
  colframe = yellow!20!black,
  colback = yellow!50,
  rightrule = 1pt,
  leftrule = 1pt,
  bottomrule = 2pt,
  colbacktitle = yellow!20!black,
  theorem style = standard,
  breakable,
  arc = 0pt
}{remark-ref}
\newtcbtheorem[number within = chapter]{myex}{例題}{
  fonttitle = \gtfamily\sffamily\bfseries\upshape,
  colframe = black,
  colback = white,
  rightrule = 1pt,
  leftrule = 1pt,
  bottomrule = 2pt,
  colbacktitle = black,
  theorem style = standard,
  breakable,
  arc = 0pt
}{example-ref}
\newtcbtheorem[number within = chapter]{exc}{Requirement}{myexc}{exc-ref}
\newcommand{\rqref}[1]{{\bfseries\sffamily 練習問題 \ref{exc-ref:#1}}}
\newtcbtheorem[number within = chapter]{definition}{Definition}{mydef}{definition-ref}
\newcommand{\dfref}[1]{{\bfseries\sffamily 定義 \ref{definition-ref:#1}}}
\newtcbtheorem[number within = chapter]{prop}{命題}{myprop}{proposition-ref}
\newcommand{\prref}[1]{{\bfseries\sffamily 命題 \ref{proposition-ref:#1}}}
\newtcbtheorem[number within = chapter]{rem}{注意}{myrem}{remark-ref}
\newcommand{\rmref}[1]{{\bfseries\sffamily 注意 \ref{remark-ref:#1}}}
\newtcbtheorem[number within = chapter]{ex}{例題}{myex}{example-ref}
\newcommand{\exref}[1]{{\bfseries\sffamily 例題 \ref{example-ref:#1}}}
% --- 再定義コマンド ---
% \mathtoolsset{showonlyrefs=true} % 必要な式番号のみ表示
\pagestyle{fancy} % ヘッダー・フッターのスタイル設定
\chead{応用量子物性講義ノート} % 中央ヘッダー
% \rhead{}
\fancyhead[R]{\rightmark}
\renewcommand{\sectionmark}[1]{\markright{\thesection\ #1}}
\cfoot{\thepage} % 中央フッターにページ番号
\lhead{}
\rfoot{Yuto Masuda and Haruki Aoki} % 右フッターに名前
\setcounter{tocdepth}{4} % 目次の深さ
\makeatletter
\@addtoreset{equation}{section} % サブセクションごとに式番号をリセット
\makeatother

% --- メタ情報 ---
\title{応用量子物性講義ノート}
\date{更新日\today}
\author{Yuto Masuda and Haruki Aoki}

\begin{document}
  \begin{myex}{バンドギャップ}{}
    長さ$L$の無限井戸型ポテンシャル中の1次元自由電子の波動関数$\phi_k(x)$とエネルギー固有値$\epsilon_0(k)$は，
    \begin{align}
      \phi_k(x) &= \braket{x}{k} = \frac{1}{\sqrt{L}}\e^{\i kx} \\ 
      \epsilon_0(k) &= \frac{\hbar^2k^2}{2m}\ \qty(k = \frac{2\pi}{L}N,\ N\in\mathbb{N})
    \end{align}
    である．
    無限井戸型ポテンシャルに$V(x + a) = V(x)$を満たすポテンシャル$V(x)$が加わったときを考える．
    $V(x)$は，
    \begin{align}
      V(x)& = 2V\cos\qty(\frac{2\pi}{a}x)\\
      & = V(\e^{\i\frac{2\pi}{a}x} + \e^{-\i\frac{2\pi}{a}x})\\
      & = V(\e^{\i gx} + \e^{-\i gx})
    \end{align}
    と書けるとする．
    ただし$g$は，
    \begin{align}
      g\coloneqq \frac{2\pi}{a}\label{wave-number-g-def}
    \end{align}
    である．
    \par
    以下の問いに答えよ．
    \begin{enumerate}
      \item 結晶中の周期ポテンシャルによりバンドギャップができることを2次までの摂動論を用いて説明せよ．
      \item 縮退のある場合の摂動論を用いてバンドギャップエネルギーを見積もれ．また，Brillouinゾーン端近傍で近似した波動関数を表せ．
      \item バンドギャップとBraggの回折条件との関係について議論せよ．
    \end{enumerate}
    \tcblower
    \begin{enumerate}
      \item バンドギャップの成り立ち \par
        2次の摂動によるエネルギー補正は，
        \begin{align}
          E_n = E_n^{(0)} + \mel**{n^{(0)}}{\hat{V}}{n^{(0)}} + \sum_{m\neq n}\frac{\abs{\mel**{m^{(0)}}{\hat{V}}{n^{(0)}}}^2}{E_n^{(0)} - E_m^{(0)}} \label{2nd-order-perturbation-eigenvalue-re}
        \end{align}
        と書ける．\refe{2nd-order-perturbation-eigenvalue-re}に対して離散Fourier変換を行うことで，\refe{2nd-order-perturbation-eigenvalue-re}の状態と
        エネルギーのラベリングを$n$から$k$に変更する．
        $V_{k'k} \coloneqq \mel**{k'}{\hat{V}}{k}$とすると，
        \begin{align}
          E(k) = \epsilon^{(0)}(k) + V_{kk} + \sum_{k'\neq k}\frac{\abs{V_{k'k}}^2}{\epsilon^{(0)}(k) - \epsilon^{(0)}(k')}
        \end{align}
        と書ける．
        摂動によるエネルギーは．
        \begin{align}
          V_{k'k} &= \frac{V}{L}\int_{-L/2}^{L/2}\phi_{k'}^{*}(x)\hat{V}(x)\phi_k(x)\dd{x} \\
          &= V\qty[\frac{\sin\qty(\frac{q_+L}{2})}{\frac{q_+L}{2}} + \frac{\sin\qty(\frac{q_-L}{2})}{\frac{q_-L}{2}}] 
        \end{align}
        と計算される．ただし$q_+$と$q_-$を，
        \begin{align}
          q_+ &\coloneqq -k' + g + k \\ 
          q_- &\coloneqq -k' - g + k
        \end{align}
        と定義した．摂動によるエネルギーは$\r{sinc}$関数の形になっているので，$L\to\infty$では規格化されたデルタ関数$\tilde{\delta}(x_1, x_2)$と解釈できる．
        よって，
        \begin{align}
          V_{k'k} = V\qty(\tilde{\delta}(q_+, 0) + \tilde{\delta}(q_-, 0)) = 
          \begin{dcases}
            V & q_+ = 0\ \r{or}\ q_- = 0\\
            0 & \r{otherwise}
          \end{dcases}
        \end{align}
        である．なお，$q_+ = q_- = 0$となるのは$g = 0$であるが，$g$の定義である\refe{wave-number-g-def}よりありえない．エネルギーは，
        \begin{align}
          E(k) = \epsilon^{(0)}(k) + \frac{V^2}{\epsilon(k) - \epsilon(k + g)} + \frac{V^2}{\epsilon(k) - \epsilon(k - g)}\label{2nd-order-perturbation-eigenvalue-in-electron}
        \end{align}
        となる．\refe{2nd-order-perturbation-eigenvalue-in-electron}の振る舞いを第1 Brillouinゾーンの内側と外側で確認する．
        ポテンシャルの対称性から右側のみを計算すればよい．
        \begin{enumerate}
          \item 第1 Brillouinゾーン内側$\qty(k = k_1 < \frac{\pi}{a})$の振る舞い\par
          $\epsilon(k)$は放物線なので，
          \begin{align}
            \begin{dcases}
              \epsilon(k_1) \ll \epsilon^{(0)}(k_1 + g) \\
              \epsilon(k_1) < \epsilon^{(0)}(k_1 - g)
            \end{dcases}
          \end{align}
          が成りたつ．よって，$E(k) < \epsilon^{(0)}(k)$が成り立ち，摂動が加わった後のエネルギーは加わる前のエネルギーより小さくなる．
          \item 第1 Brillouinゾーン外側$\qty(k = k_2 > \frac{\pi}{a})$の振る舞い\par
            第1 Brillouinゾーン内側のときと同様に考えると，
            \begin{align}
              \begin{dcases}
                \epsilon(k_2) \ll \epsilon^{(0)}(k_2 + g)\\
                \epsilon(k_2) > \epsilon^{(0)}(k_2 - g)
              \end{dcases}
            \end{align}
            を得る．よって，$E(k) > \epsilon^{(0)}(k)$が成り立ち，摂動が加わった後のエネルギーは加わる前のエネルギーより大きくなる．
          \end{enumerate}
          以上の議論により，結晶中の周期ポテンシャルによりバンドギャップが形成されることがわかった．
      \item バンドギャップエネルギーの見積もりと波動関数\par
        \refe{2nd-order-perturbation-eigenvalue-in-electron}に$k = \pm\frac{\pi}{a}$を代入すると発散してしまう．
        以下では2重縮退があるときの摂動を考えバンドギャップエネルギー$\Delta E$を求める．$k_+$と$k_-$を，
        \begin{align}
          k_+ &\coloneqq \frac{\pi}{a} \\ 
          k_- &\coloneqq -\frac{\pi}{a}
        \end{align}
        と定義する．$V_{k_+k_+} = V_{k_-k_-} = 0$，$V_{k_+k_-} = V_{k_-k_+} = V$なので，
        \refe{1st-order-energy-with-degeneracy}より，1次摂動によるエネルギーは，
        \begin{align}
          E_n^{(1)} &= \frac{1}{2}\qty[(V_{k_+k_+} + V_{k_-k_-}\pm\sqrt{(V_{k_+k_+} - V_{k_-k_-})^2 + 4\abs{V_{k_+k_-}}^2})] \\ 
          &= \pm V
        \end{align}
        である．2つのエネルギー補正の差がバンドギャップエネルギー$\Delta E$と解釈できるので，
        \begin{align}
          \Delta E = 2V
        \end{align}
        を得る．\par
        \refe{degeneracy-n0-def}の$\alpha$と$\beta$は\refe{secular-eq}の解であるから，$\sqrt{\alpha^2 + \beta^2} = 1$なる規格化条件を課すと，
        \begin{align}
          \mqty(\alpha \\ \beta) = 
          \begin{dcases}
            \frac{1}{\sqrt{2}}\mqty(1 \\ 1) & E_n^{(1)} = V \\ 
            \frac{1}{\sqrt{2}}\mqty(1 \\ -1) & E_n^{(1)} = -V
          \end{dcases}
        \end{align}
        となる．よって，
        $\Delta E = \pm V$に対応する波動関数はBrillouinゾーン端で，
        \begin{align}
          \psi_{+}& = \frac{1}{\sqrt{2}}\qty(\phi_{\pi/a} + \phi_{-\pi/a})\propto\cos\qty(\frac{\pi}{a}x)\label{degeneracy-wave-function-positive}\\
          \psi_{-}& = \frac{1}{\sqrt{2}}\qty(\phi_{\pi/a} - \phi_{-\pi/a})\propto\sin\qty(\frac{\pi}{a}x)\label{degeneracy-wave-function-negative}
        \end{align}
        であり，定在波が生じる．
    \item バンドギャップとBragg反射\par
      バンドギャップの起源はBragg反射である．Bragg反射は，
      \begin{align}
        2a\sin\theta = \lambda
      \end{align}
      を満たす．
      今回の場合は1次元なので$\theta = \pi/2$であり，波数は$k = 2\pi/\lambda$である．よって，Bragg条件は
      \begin{align}
        k = \frac{\pi}{a}\label{1d-bragg-k-condition}
      \end{align}
      と書き換えられる．つまり，\refe{1d-bragg-k-condition}を満たす波数のみが反射し定在波をつくる．$V(x) = 2V\cos\qty(\frac{2\pi}{a}x)$であったから，
      \refe{degeneracy-wave-function-positive}で表される波動関数はポテンシャルが最小となる波数で確率振幅が最大となる．
      \refe{degeneracy-wave-function-negative}で表される波動関数はポテンシャルが最大となる波数で確率振幅が最大となる．
      よって，エネルギーは$\psi_+$が$\psi_-$より低くなる．これによりバンドギャップが生じる．
  \end{enumerate}
  \end{myex}
\end{document}
